% Source: https://artofproblemsolving.com/community/c3h1458903_comc_1997_4
% Stable identity: aops-topic-1458903-q04
None of the digits is $0$, since a number such as $\overline{0a}$ is not a two-digit number. The six numbers are $\overline{ab},\overline{ba},\overline{ac},\overline{ca},\overline{bc},\overline{cb}$. Each digit is a tens digit in two of them and a units digit in two of them, so the sum is
\[
2\cdot10(a+b+c)+2(a+b+c)=22(a+b+c).
\]
So $22(a+b+c)=484$, i.e.\ $a+b+c=22$. Three distinct digits with sum $22$ must include $9$ (otherwise the sum is at most $8+7+6=21$); the other two then sum to $13$ and are distinct digits other than $9$: $\{5,8\}$ or $\{6,7\}$ ($\{4,9\}$ would repeat the $9$). The sets are
\[
\boxed{\{5,8,9\}\ \text{ and }\ \{6,7,9\}}.
\]
